\chapter*{Заключение} \label{ch-conclusion}
\addcontentsline{toc}{chapter}{Заключение}

В работе выполнен сравнительный анализ фреймворков модульного тестирования JavaScript: Jest, Vitest и Mocha. Для сравнения разработано компактное программное средство на Node.js, которое импортирует данные из JSON и CSV, запускает одинаковые модульные тесты и фиксирует время выполнения, пиковое потребление оперативной памяти и пиковую загрузку CPU.

На авторском JSON-наборе из 5000 элементов медианное время составило 10640.290 мс для Jest, 16273.677 мс для Vitest и 3606.863 мс для Mocha. Медианная пиковая память составила 90.781 МБ, 372.242 МБ и 57.340 МБ соответственно. На CSV-наборе из 15 элементов общая тенденция сохранилась: Mocha оказался самым быстрым и самым легким по памяти, Vitest показал наибольшую пиковую память, Jest занял промежуточное положение.

Главный вывод состоит в том, что выбор фреймворка должен учитывать контекст. Mocha подходит для компактных проектов и случаев, где важна минимальная инфраструктурная нагрузка. Jest удобен для проектов, которым нужны развитые встроенные средства тестирования. Vitest рационален для проектов, связанных с Vite и современной ESM-сборкой, но в данном эксперименте его инфраструктура потребовала больше памяти.

Все обязательные элементы задания выполнены: есть математическая постановка, иллюстративный пример, псевдокод в окружении \texttt{algorithm}, пошаговое выполнение метода, описание программного обеспечения, инструкция запуска, эталонные и авторские данные, \texttt{requirements.txt} и автоматически формируемый список источников из BibTeX.
